% !TeX program = xelatex
% ============================================================================
%  第 1 课　先导课：抽象代数是什么　—— 教师版讲义
%  与 02-课件/第01课-先导课/第01课-先导课.tex 同步（以课件为准，改课件请回来同步这里）
% ============================================================================

\jhlesson{1}{先导课：抽象代数是什么}{时长 45 分钟\quad·\quad 教具：无\quad·\quad 本课不发单页}

\begin{mubiao}
  \begin{itemize}
    \item 知道抽象代数研究的是「运算本身的规律」，不是一批新公式。
    \item 能用集合的说法描述一堆东西（只会写、会读，不做集合运算）。
    \item 记住一个悬念：999999 能被 7 整除，为什么？
  \end{itemize}
\end{mubiao}

\noindent{\small 本课节奏：0—8 一个吓人的名字｜8—18 两个小例子｜18—26 集合初步｜26—36 一个悬念｜36—45 课程地图。全程至少 4 次「用手指回答」。}

\section*{一、一个吓人的名字（0—8 分钟）}

\noindent 板书四个字：\textbf{抽象代数}。

先给学生看课程名《从数数开始的抽象代数生活》——它化用了动画《Re:从零开始的异世界生活》，
但很遗憾这门课和 Re0 没有任何关系。如果不起一个有梗的名字，这门课只有四个字：抽象代数。
（课件第 2 页放了动画海报，第 5 页会给他们看一眼大学教材的目录。）

\begin{definition}[抽象代数]
  它是大学数学系的专业课，一般大二才开。它研究的不是「数怎么算」，而是「算」本身。
\end{definition}

\subsection*{三个台阶}

先回到学生能看懂的语言：数字的世界是怎么一点一点变抽象的？

在黑板上竖着写三行，一边写一边说：

\begin{itemize}
  \item \textbf{小学}：只用数字。\quad \(3+5=8\)
  \item \textbf{初中}：开始用字母代替数。\quad \(a+b=b+a\)
  \item \textbf{抽象代数}：连「运算」本身也用字母代替。
\end{itemize}

第三级台阶的意思是：不管是数的加法、钟表上的走针，还是把队伍的次序换一下，
只要它们遵守同样的规矩，数学家就把它们归到一类里研究。

\begin{zhu}
  这句话是本课的题眼，说慢一点。学生此刻不需要听懂，只需要感到「好像有点道理」。
\end{zhu}

\subsection*{这门课有两个约定}

\begin{itemize}
  \item 我们尽量少记忆一些术语和定义；
  \item 只看生活里的例子：钟表、体育课、三角形、正方形、身份证。
\end{itemize}

最后一节课，你要自己证明一个二百多年前的结论。

\begin{caozuo}
  这里可以停一下，卖个关子：「你们现在肯定不相信。没关系，我们先把话放在这儿。」
\end{caozuo}

\section*{二、生活里的两个例子（8—18 分钟）}

\begin{example}[例子一：体育课上的队形]
  \textbf{问题}：体育课上，老师要把两列纵队变成四列纵队，他做了什么？
\end{example}

先把问题念出来，让学生想 10 秒，然后请他们用手指表示「我在原队列里该往哪挪一步」：
往左挪的举 \textbf{1} 根，往右挪的举 \textbf{2} 根。

一句总结：\emph{人一个没变，动的只是位置。}

\begin{example}[例子二：半夜 12 点]
  \textbf{问题}：为什么半夜 12 点，既叫 0 点，又叫 24 点？
\end{example}

在黑板左侧画一个圆，写上 \(0,3,6,9,12,15,18,21,24\)，从 0 开始一段一段往右描，一直描到 24。
描到 24 时停下，问：这个位置刚才在哪儿？

（课件第 8 页放了钟面图。）

\noindent 手指回答：老师问「早上 7 点和晚上 7 点，在钟面上是同一个位置吗？」
认为是的举 1 根，不是的举 2 根。

\noindent 把两个例子并起来板书：

\begin{center}
\begin{tabular}{@{}ll@{}}
\toprule
队形 & 人没变，位置上重新排了一遍。 \\
时间 & 表没变，指针绕了一圈。 \\
\midrule
\textbf{共同点} & \textbf{东西没变，只是「重新排了一下」。} \\
\bottomrule
\end{tabular}
\end{center}

这门课要做的，就是把这句话说清楚。

\section*{三、集合初步（18—26 分钟）}

\begin{definition}[集合]
  把确定的一堆东西放在一起，就叫做一个\textbf{集合}，用大括号写出来。
\end{definition}

\begin{example}
  \(\{1,2,3\}\)；\quad \{甲, 乙, 丙\}；\quad \{红, 黄, 蓝\}。

  所有的有理数放在一起也是集合；自然数、整数也都可以看作集合。
\end{example}

\begin{zhu}
  今天只需要知道「怎么写、怎么读」，不涉及集合的运算（交、并、补都不讲）。
\end{zhu}

然后立刻给出这节课马上要用的三个写法：

\begin{itemize}
  \item \(\{1,2,3,\dots,12\}\) —— 钟面上的 12 个时刻；
  \item \(\{A,B,C\}\) —— 三角形的三个顶点；
  \item \{原地不动, 转 \(90^\circ\), 转 \(180^\circ\), 转 \(270^\circ\)\} —— 正方形的四种转法。
\end{itemize}

要停一下的地方：指着最后一行说，\emph{注意，这个大括号里装的不是数，是四个动作。}
稍作停顿，不用多解释，学生能接受「一堆动作也能打包」就可以往下走。

\noindent 手指回答：老师问「\(\{\text{转 }90^\circ,\ \text{转 }180^\circ\}\) 这样的写法，符合刚才的定义吗？」
认为符合的举 1 根，不符合的举 2 根。答案是符合——集合只管「确定的一堆东西」，不挑它们是数还是动作。

\section*{四、一个悬念（26—36 分钟）}

转身在黑板中间写下：

\[
  999999 \div 7 = \underline{\hspace{5em}}
\]

先让全班猜能不能整除，再当场算：\(999999 \div 7 = 142857\)，正好整除，一点不剩。
停两秒，然后问：\emph{为什么它就能被 7 整除？}——今天不回答。

给一条线索（这条线索也是最后一节课的钥匙）：

\[
  999999 = 1000000 - 1 = 10^{6} - 1
\]

于是问题变成：为什么 \(10^{6}-1\) 能被 7 整除？这件事和素数 7 有什么关系？把 7 换成别的数还管用吗？

\begin{caozuo}
  把 \(999999 \div 7\) 这一行写在黑板\textbf{最右边}，旁边标注「第 13 课回来回答」。
  这个问号整门课都不擦——第 1 课写上去，第 13 课当场划掉。
\end{caozuo}

顺手再玩一个把戏，把学生的好奇心拉满。课件分两步翻，先给前三行，让学生自己猜后三行：

\begin{example}[142857 的循环]
  \begin{center}
  \begin{tabular}{@{}ll@{\hspace{2.5em}}ll@{}}
    \(142857 \times 1 = 142857\) & & \(142857 \times 4 = 571428\) & \\
    \(142857 \times 2 = 285714\) & & \(142857 \times 5 = 714285\) & \\
    \(142857 \times 3 = 428571\) & & \(142857 \times 6 = 857142\) & \\
  \end{tabular}
  \end{center}
\end{example}

数字还是那六个，只是换了个位置——今天第二个例子里说的「重新排一下」，又出现了。

\noindent 手指回答：老师问「这六个结果里，是不是每个结果都只用到了 1、4、2、8、5、7 这六个数字？」
认为是的举 1 根，不是的举 2 根。

\section*{五、这门课要去哪（36—45 分钟）}

板书课程地图——四站路：

\begin{itemize}
  \item \textbf{第一站　钟表与取余}：数着数着，会绕回来。
  \item \textbf{第二站　三角形与正方形}：位置换一换，还是那几个样子。
  \item \textbf{第三站　换位置的算术}：先做这个再做那个，结果会不一样。
  \item \textbf{第四站　费马小定理}：回到今天那个 999999。
\end{itemize}

最后一节课，你会带走三样东西：

\begin{itemize}
  \item 一个能自己证明的结论（费马小定理）；
  \item 一套看世界的办法：把不同事情里相同的规矩挑出来；
  \item 一个终于能回答的问题：999999 为什么能被 7 整除。
\end{itemize}

\noindent 收尾一句话：\emph{这门课的终点，是你能自己证明黑板右边那个问题。}
说完指回黑板最右边那行字，下课前不要擦掉。

\subsection*{如果你想读一点参考书}

这些书你大概率是读不懂的，至少也得上了高中（课件最后一页有完整书目，这里只列两条最容易找的）：

\begin{itemize}
  \item 人民教育出版社．普通高中课程标准实验教科书 数学 选修 3-4 A 版：对称与群．北京：人民教育出版社，2007．
  \item 内森·卡特．群论：彩图版．郭小强，罗翠玲，译．北京：机械工业出版社，2019．
\end{itemize}

\begin{xiaojie}
  \begin{itemize}
    \item 抽象代数 = 研究「运算这件事本身」的规律，不背新公式。
    \item 集合 = 把确定的一堆东西用大括号打包；动作也可以打包。
    \item 今天的两个例子说的是同一句话：\textbf{东西没变，只是重新排了一下}。
    \item 悬而未决：\(999999\) 为什么能被 7 整除？（第 13 课回答）
  \end{itemize}
\end{xiaojie}

\noindent\textbf{下节课}：从余数到模运算——我们会给今天那个「绕圈」装上记号，并第一次见到单位元和逆元。